% TOP_PROBLEM: {"id": 3000003, "title": "Acyclic orientation with parity constraints", "category_id": 2, "category": {"id": 2, "name": "combinatorics", "display_name": "Combinatorics", "description": "Counting problems, graph theory, discrete structures.", "slug": "combinatorics", "order_index": 2, "created_at": "2026-07-31T15:26:25.671Z"}, "status": "partially_solved", "problem_number": "AMR-029-0003", "set_id": 13}
% FIRST_POSTED: 2026-10-02
% ENTRY_KIND: statement_only
% UnsolvedMath ID 3000003; source catalogue code AMR-029-0003
% !TeX program = lualatex
% Definition and sources only; this file is not an LLM solution attempt.
\documentclass[11pt,a4paper]{article}
\usepackage[margin=25mm]{geometry}
\usepackage{fontspec}
\setmainfont{lmroman10-regular.otf}[BoldFont=lmroman10-bold.otf,
 ItalicFont=lmroman10-italic.otf,BoldItalicFont=lmroman10-bolditalic.otf]
\usepackage{amsmath,amssymb,mathtools}
\usepackage{unicode-math}
\setmathfont{latinmodern-math.otf}
\AtBeginDocument{\let\setminus\smallsetminus}
\usepackage{xurl,hyperref}
\hypersetup{colorlinks=true,urlcolor=blue}
\setcounter{secnumdepth}{0}
\setlength{\parindent}{0pt}
\setlength{\parskip}{5pt}
\setlength{\emergencystretch}{3em}
\begin{document}
\section*{Acyclic orientation with parity constraints}\label{3000003}
{\small UnsolvedMath ID 3000003 · AMR-029-0003 · Combinatorics · AMR Open Problem Lists}
\subsection{Definitions and mathematical statement}
Let $G=(V,E)$ be a finite loopless undirected graph; parallel edges,
if present, are counted separately. An orientation chooses one direction
for each edge. It is acyclic if it contains no directed cycle.
For $T\subseteq V$, a $T$-parity orientation satisfies
\[
 d_D^-(v)\equiv\begin{cases}1&v\in T,\\0&v\notin T\end{cases}
                       \pmod2,
\]
where $d_D^-(v)$ counts arcs entering $v$.

Give a structural, effectively usable characterization of the pairs
$(G,T)$ admitting an acyclic $T$-parity orientation. The all-even case
$T=\varnothing$ is a separately highlighted subproblem. An equivalent
ordering formulation asks for a linear ordering of $V$ in which each
vertex has an odd number of incident edges to earlier vertices exactly
when it belongs to $T$.

The third source question asks which graphs have this property for
every proper subset $T\subsetneq V$ satisfying
\[
                         |T|\equiv|E|\pmod2.
\]
An orientation may be chosen separately for each such $T$.
The parity condition is necessary because the sum of all indegrees is
$|E|$; every connected component also has its own parity constraint.
Excluding $T=V$ is natural because a nonempty finite acyclic digraph
has a source of indegree zero. These necessary observations alone are
not offered as a sufficient characterization. Exhaustive enumeration
of vertex orders is finite, but does not provide the sought structural
description or an efficient recognition method.
\subsection{Short English statement}
Characterize acyclic graph orientations with prescribed indegree parity, including graphs supporting every admissible proper parity set.
\subsection{Sources}
\begin{itemize}
\item[S1] ULAM.AI and the UnsolvedMath Contributors, supplied problems.json, record 3000003 (AMR-029-0003), AMR Open Problem Lists. Catalogue identity and source transcription; CC BY 4.0.
\url{https://huggingface.co/datasets/ulamai/UnsolvedMath}.
\item[S2] Egres Open - List of open problems, Open-problem page: Acyclic orientation with parity constraints. Original source indexed by AMR-029-0003.
\url{https://oldlemon.cs.elte.hu/egres/open/List_of_open_problems}.
\item[S3] EGRES Open, Acyclic orientation with parity constraints. Individual source statement and remarks.
\url{https://lemon.cs.elte.hu/egres/open/Acyclic_orientation_with_parity_constraints}.
\end{itemize}
\end{document}
